\section{Separated exponential scales}\label{se:sec:separated}
\begin{theorem}\label{se:main}
For \(p=2,3,\ k\ge1\), the construction has support \(\splitcount{p+2}{k+1}\) and at least \(8k\lfloor(9p-2)/4\rfloor\) hyperbolic limit cycles. We use \((p,k)=(2,1),(3,4)\) at six and ten terms. Write \(L=\Mclr\) for the even exponent-clearing integer. The actual degree \(E_N\) is the maximum of \(1+L(3+\rho+\alpha_0)\), \(1+L(4+\rho)\), \(1+L(3+q_j+\alpha_j)\) for \(1\le j\le p\), and \(1+3L\). Requiring \(0<\delta<1/16\) at the small-scale step makes the clearing factor exceed 16.
\end{theorem}
Section~\ref{sup:primitive-fields} gives the full common factor of the six- through ten-term realizations and the actual degree after cancellation.

\subsection{One center with several energy scales}
Fix \(k\) and put \(b_i=\delta^i\), \(1\le i\le k\), for \(0<\delta<1\).
Set \(B=\sum_{i=1}^k b_i\) and
\begin{equation}
 W_\delta(u)=B(e^u-1)+\sum_{i=1}^k(1-e^{b_i u}).                    \label{se:W}
\end{equation}
It has a nondegenerate minimum \(W_\delta(0)=0\).
Indeed \(W_\delta'=\sum_i b_i(e^u-e^{b_i u})\) has the sign of \(u\)
and \(W_\delta''(0)=\sum_i b_i(1-b_i)>0\). The potential tends to
\(k-B\) on the left and to infinity on the right. For
\(0<E<k-B\), \(W_\delta(u)+Z^2=E\) is one regular oval surrounding the
origin.

For \(A(w)=\sum_j c_jw^j\) let
\[
 {\cal I}_{\delta,A}(E)=\iint_{W_\delta(u)+Z^2<E} A(Z^2)\dd u\dd Z .
\]
On \(0<h<1\) define
\begin{equation}
 J_j(h)=\frac{2}{2j+1}\int_0^h\frac{(h-s)^{j+1/2}}{1-s}\dd s,
 \qquad J[A](h)=\sum_j c_jJ_j(h).                                   \label{se:mom}
\end{equation}
For every fixed \(h\) in this interval,
\begin{equation}
 b_i\,{\cal I}_{\delta,A}(i-1+h)\longrightarrow J[A](h)
 \quad(\delta\downarrow0),\qquad i=1,\ldots,k.                       \label{se:decouple}
\end{equation}
To prove this, put \(x=-b_i u\). At fixed \(x>0\), the potential tends
to \(i-1+1-e^{-x}\). The earlier exponentials tend to zero and the
later ones tend to one. At \(x<0\) the term \(B e^{|x|/b_i}\) dominates
all the others, since \(B\ge\delta\) and \(b_j\le\delta\). It tends to
infinity even after division by \(e^{\delta|x|/b_i}\).
For any fixed \(R>-\log(1-h)\), the entire sublevel set in \((x,Z)\)
is eventually contained in \([-1,R]\times[-\sqrt{k},\sqrt{k}]\), by
monotonicity on each side. Dominated convergence and the substitution
\(s=1-e^{-x}\) give \eqref{se:decouple}. In particular, finitely many strict
signs at fixed regular energies persist simultaneously at all \(k\) scales.

\subsection{Moments at the two ends of an interval}
As \(h\downarrow0\),
\[
 J_j(h)\sim\frac{4}{(2j+1)(2j+3)}h^{j+3/2}.
\]
For the upper endpoint put \(v=1-h>0\). Polynomial division in the
variable \(y=h-s\) gives
\begin{equation}
 \frac{2j+1}{2}J_j(h)=
 \sum_{l=0}^{j}(-v)^l\frac{h^{j-l+1/2}}{j-l+1/2}
      +2(-1)^{j+1}v^{j+1/2}\arccos\sqrt v .                         \label{se:explicit}
\end{equation}
Thus \(J_j(1-v)\) is analytic in \(\sqrt v\) at zero and its first
noninteger power is
\((-1)^{j+1}2\pi v^{j+1/2}/(2j+1)\).
Writing \(y_j=j+1/2\), its finite Taylor data in \(h-1\) are
\begin{equation}
 T_{lj}=\frac{J_j^{(l)}(1^-)}{l!}
   =\frac{\binom{y_j}{l}}{y_j(y_j-l)},\qquad 0\le l\le j .          \label{se:matrix}
\end{equation}
This follows by differentiating \eqref{se:mom} and taking the endpoint
limit of the integrable power \((1-s)^{j-l-1/2}\).

\begin{lemma}\label{se:det}
For \(a\ge0\), take rows \(0,\ldots,a\) of \eqref{se:matrix} and any
\(a+1\) distinct integer columns \(j\ge a\). The resulting matrix is
nonsingular.
\end{lemma}
\begin{proof}
Remove nonzero column and row factors and put \(Y=2j+1\).
The row functions become
\(\prod_{b=0}^{l-1}(Y-2b)/(Y-2l)\).
Multiply each column by \(\prod_{b=0}^a(Y-2b)\).
The resulting row polynomial \(P_l(Y)\) reduces modulo two to
\(Y^{a+l}\). The integer polynomial
\(\det(P_l(Y_s))_{l,s=0}^a\) is alternating and is divisible in
\(\mathbb Z[Y_0,\ldots,Y_a]\) by
\(\Delta=\prod_{i<j}(Y_j-Y_i)\).
Modulo two the determinant is exactly
\(\Delta\prod_sY_s^a\). Cancelling the nonzero polynomial \(\Delta\)
over \(\mathbb F_2\) shows that the integer quotient is odd at any tuple
of odd integers. At the distinct odd nodes of the lemma the Vandermonde
is nonzero as well.
\end{proof}

Define \(A_{2j}\) on the powers \(w^j,\ldots,w^{2j}\) so its data in
\eqref{se:matrix} through order \(j\) are \(0,\ldots,0,1\).
Define \(A_{2j+1}\) by taking coefficient one on \(w^j\) and solving on
\(w^{j+1},\ldots,w^{2j+1}\) to kill the same \(j+1\) data.
Lemma~\ref{se:det} gives rational solutions. For odd indices, omission of
the highest power would give a nonzero vector in the kernel of the even
system. For even indices \(2j>0\), omission of the last power would give
a kernel on the first \(j\) rows and the \(j\) remaining columns. Thus
each degree is exactly its index, including \(A_0=1/4\).

\begin{lemma}\label{se:zeros}
For every \(d\ge1\), there is a rational polynomial \(A\) of degree
\(d\) for which \(J[A]\) has at least \(2\lfloor(3d+1)/4\rfloor\)
changes of sign on disjoint regular open intervals in \((0,1)\).
\end{lemma}
\begin{proof}
Write \(e_l=\lceil l/2\rceil\). For \(r=\sqrt{1-h}\) the leading term
of \(J[A_l](1-r^2)\) is \((-1)^{e_l}c_l r^l\), with \(c_{2j}=1\) and
\(c_{2j+1}=2\pi/(2j+1)>0\), by \eqref{se:explicit}.
For a sufficiently small positive rational \(q\) use
\begin{equation}
 A=\sum_{l=0}^d(-1)^{l+e_l}q^{(d+1-l)^2-1} A_l .                  \label{se:density}
\end{equation}
At \(r=q^{2(d+1-l)}\) the direction of index \(l\) dominates. Indeed
the leading exponent of \(q\) at index \(v\) exceeds that at \(l\) by
\((v-l)^2\). The \(d+1\) signs near \(h=1\) are therefore \((-1)^l\).

At \(h=q^{4(d-2j)}\), \(0\le j\le\lfloor(d-1)/2\rfloor\), the term of
index \(2j+1\) dominates near zero. Divide out \(h^{3/2}\). The odd
term \(2i+1\) starts with a positive multiple of \(h^i\), and its leading
exponent of \(q\) exceeds the chosen one by \(4(i-j)^2\). Each even term
\(2i\) starts at power at least \(i\) and costs an additional
\(2(d-2i)+1>0\) compared to the odd exponent formula at \(i\). This
comparison also covers \(2i=d\) by algebraic extrapolation of that
formula. The signs at the lower marks in increasing \(h\) are
\((-1)^j\).

The upper list gives \(d\) changes and the lower list gives
\(\lceil d/2\rceil-1\). In the concatenated list the first and the
last signs are both positive. Its number of changes is even, hence
at least the smallest even integer at least \(d+\lceil d/2\rceil-1\),
which is \(2\lfloor(3d+1)/4\rfloor\). For small \(q\) the lists are
disjoint and all marks have rational energies.
\end{proof}

Choose \(A\) as in \eqref{se:density}, then choose a sufficiently small
rational \(\delta>0\) to preserve the signs by \eqref{se:decouple} and put
all translated marks below \(k-B\). On every compact regular interval in
\((0,k-B)\) the functions \({\cal I}_{\delta,A}\) and
\({\cal I}_{\delta,1}>0\) are real analytic. For example
\(v=\operatorname{sgn}(u)\sqrt{W_\delta(u)}\) has an increasing analytic
inverse at and between the turning points. Integration over
\(v^2+Z^2<E\) then has an analytic integrand on a rescaled fixed disk.
The ratio \({\cal I}_{\delta,A}/{\cal I}_{\delta,1}\) is nonconstant on
each closed sign-changing bracket, and its derivative has finitely many
zeros on their union. Add a single small rational constant \(c\) to
\(A\), preserving all strict signs and avoiding the negatives of these
critical values. The resulting integral has a simple zero in every
selected bracket. This applies to all scales of the same potential.

\subsection{Realizing the density with \texorpdfstring{\(p+2\)}{p+2} powers}
We give the argument with general \(p\ge2\) since dropping the horizontal
tilt changes the rank condition from the one used for the paired wells.
Set \(d=3p-1\), \(N_j=2j\) for \(0\le j<p\), and \(N_p=2p-1\).
Let \(v_p=1\) and choose the integers \(v_0,\ldots,v_{p-1}\) making
\(B_0(z)=\sum_{j=0}^p v_j(1+z)^{N_j}\) odd. The triangular even-power
system fixes them, and gives \((2,-3,1)\) for \(p=2\) and \((-16,20,-5,1)\) for \(p=3\).
Here \(\ev\) denotes the even part in the displayed variable. Write
\(L=1+z\) and define
\[
 h_j^V=-z^{-1}\int_0^z((1+y)^{N_j}-1)\dd y,\qquad
 h_*^V=z^{-1}\int_0^z y\log(1+y)\dd y.
\]
Take rows \([z^{2i}],\ 0\le i\le d\), on the following functions, in order,
to form a \(3p\) by \(3p\) rational matrix \(J\).
\begin{gather}
 \ev L^{N_j}\ (0\le j<p),\qquad
 \ev(L^{N_j}\log L)\ (1\le j\le p),\notag\\
 -\ev(B_0h_j^V)'\ (1\le j<p),\qquad
 \ev\big(v_0\log L-(B_0h_*^V)'\big).                              \label{se:vertical}
\end{gather}

\begin{lemma}\label{se:rankv}
The matrix \(J\) is nonsingular for \(p=2,3\).
\end{lemma}
\begin{proof}
For the columns and rows in the displayed order, rational elimination yields \(\det J=-1/4762800\) and \(-4/35153326333125\), respectively. These fractions can be checked by expanding \(\log(1+z)=\sum_{l=1}^{2d+2}(-1)^{l+1}z^l/l+O(z^{2d+3})\) before the indicated integrations and differentiations. The accompanying program recomputes both determinants.
\end{proof}

Write the chosen, regularly shifted density as \(A(w)=\sum_{i=0}^d c_i w^i\).
For small positive rational \(t\) take \(\lambda=t^{2d+2}\) and put
\begin{equation}
 (a_0,\ldots,a_{p-1},h_1,\ldots,h_p,s_1,\ldots,s_{p-1},\tau)^{\mathsf T}
 = J^{-1}(c_i t^{2(d-i)})_{i=0}^d .                                \label{se:controls}
\end{equation}
Fix \(a_p=s_p=0\) and an integer
\(K>2+\|J^{-1}\|_\infty\max_i|c_i|\), using the maximum row sum norm.
Set
\begin{align*}
 \gamma_j&=v_j+\lambda a_j\quad(0\le j\le p),&
 C_j&=\lambda(s_j+Kv_j)\quad(1\le j\le p),\\
 q_j&=N_j+\lambda h_j/v_j\quad(1\le j\le p),&
 C_0&=1-\textstyle\sum_{j=1}^p C_j,\qquad \rho=\lambda\tau .
\end{align*}
In the log plane put \(U=e^u,\ L=e^v=1+z\) and use
\begin{align}
 \dot u&=C_0e^{\varepsilon\gamma_0 u/C_0}L^{1+\rho}-L^{2+\rho}
                +\sum_{j=1}^p C_j e^{\varepsilon\gamma_j u/C_j}L^{1+q_j},
                \notag\\
 \dot v&=\tfrac12 t^2 \left(BU-\sum_{i=1}^k b_i U^{b_i}\right).       \label{se:logfield}
\end{align}
At \(\varepsilon=0\) the field is \((-H_v,H_u)\), with
\[
 H=\tfrac12 t^2W_\delta(u)+V(z),\quad V(0)=0,\quad
 V'(z)=zL^\rho-\sum_{j=1}^p C_j(L^{q_j}-L^\rho).
\]
Since \(V''(0)\to1\), there is a uniform analytic increasing vertical
coordinate \(Y=z\sqrt{2V(z)/z^2}\) with inverse \(z=z(Y)\).
In the coordinates \(u,Z=Y/t\) the energy is exactly
\(E=2H/t^2=W_\delta(u)+Z^2\).

In the \(\varepsilon\)-variation of the field, Green's formula on the
counterclockwise Hamiltonian oval gives the area integral of the divergence
as the first variation \(M(E)\) of its one-turn displacement in \(H\).
Explicitly the energy integral is \(\oint(\partial_\varepsilon\dot u)|_0
\dd v\), and varying the endpoint time gives zero at first order.
After \(dv=dz/L\), the density is
\(D_\lambda(z)=\gamma_0L^\rho+\sum_{j=1}^p\gamma_j L^{q_j}\).
Its even part in the energy coordinate satisfies, on bounded \(Z\)-sets,
\begin{equation}
 \frac{\ev[D_\lambda(z(Y))z'(Y)]\big|_{Y=tZ}}{\lambda t^{2d}}
                    =A(Z^2)+O(t^2).                                  \label{se:transport}
\end{equation}
For completeness, at \(\lambda=0\) the density is odd. Differentiating
at fixed controls transports it under
\(Y=z+\lambda H_1(z)+O(\lambda^2)\), where
\(H_1=\tau h_*^V+\sum_{j=1}^p(s_j+Kv_j)h_j^V\).
The first transport term is \(-\ev(B_0H_1)'\).
The part proportional to \(K\) vanishes since
\(\sum_{j=1}^p v_j h_j^V=-z^{-1}\int_0^z B_0(y)\dd y\) is odd.
The remaining Taylor coefficients to order \(2d\) are exactly
\eqref{se:vertical} applied to \eqref{se:controls}. The controls are uniformly
bounded. Analytic inversion on a fixed small disk bounds the nonlinear
remainder by \(O(\lambda^2)\) and the omitted part of the linear even
series by \(O(Y^{2d+2})\). This proves \eqref{se:transport}, also with
one derivative on compact \(Z\)-sets.

For the fixed \(\delta\), \eqref{se:transport} and the regular disk
parametrization after Lemma~\ref{se:zeros} give
\[
        M(E)/(\lambda t^{2d+1})\longrightarrow{\cal I}_{\delta,A}(E)
              \quad\text{in }C^1
\]
on compact regular energy intervals. Around each selected simple zero
take a smaller closed interval with opposite endpoint signs and derivative
of fixed nonzero sign. All these properties hold for the first variation
at one sufficiently small rational \(t\). Fix this \(t\). The one-turn
map on a transverse section with coordinate \(E\) depends analytically
on \((E,\varepsilon)\). Its displacement divided by \(\varepsilon\)
tends to \(2M(E)/t^2\) in \(C^1\). All sufficiently small positive
rational \(\varepsilon\) therefore give a simple fixed point in each
interval. The multiplier is positive and differs from one. Each cycle
stays in its own compact annular neighborhood and surrounds the
equilibrium \((0,0)\) in log coordinates. The disjoint energy intervals
give distinct nested hyperbolic cycles for the same field.

\subsection{Expansion in the real variables}
Put \(\alpha_j=\varepsilon\gamma_j/C_j\) and take \(\Mclr{}\) twice the
least common multiple of the reduced denominators of all
\(\alpha_j,\rho,q_j,b_i\), assigning denominator one to zero.
In each open quadrant make the substitution
\(u=\Mclr{}\log|x|,\ v=\Mclr{}\log|y|\) and multiply the pulled-back field by the
positive factor \(\Mclr{} x^{\Mclr} y^{\Mclr}\). The resulting polynomial has components
\begin{align}
 P={}&C_0 x^{1+\Mclr{}(1+\alpha_0)}y^{\Mclr{}(2+\rho)}
       -x^{1+\Mclr{}}y^{\Mclr{}(3+\rho)}
        +\sum_{j=1}^p C_j x^{1+\Mclr{}(1+\alpha_j)}y^{\Mclr{}(2+q_j)},\notag\\
 Q={}&\tfrac12t^2 B x^{2\Mclr{}}y^{1+\Mclr{}}
       -\tfrac12t^2\sum_{i=1}^k b_i x^{\Mclr{}(1+b_i)}y^{1+\Mclr{}}.            \label{se:poly}
\end{align}
The scaled rational exponent contributions are even integers and all
powers are nonnegative for the choices above. The limiting \(y\)-powers
in \(P\), divided by \(\Mclr{}\), are \(2,3,2+N_1,\ldots,2+N_p\), all
distinct. We have \(C_0\to1\) and \(|C_j|>2\lambda\) for \(j\ge1\).
The \(k+1\) distinct horizontal exponents come from \(1,b_1,\ldots,b_k\)
and have nonzero coefficients. Thus the exact support is \(\splitcount{p+2}{k+1}\).
The diffeomorphisms and the positive time changes preserve the section
multipliers in all four quadrants. This proves the family in
Theorem~\ref{se:main}.
